\section{Limpieza del Dataset}

Primeramente, haremos modificaciones para a los valores de las fechas (nacimiento, fecha de desaparición, fecha de registro) para pasarlos a un valor \texttt{datetime}. Esto nos ayudará a poder discretizar los valores de edad, mes y año de desaparición. En particular, con edad haremos una segunda discretización para trabajar con grupos de edad (0-12, 13-17, 18-30, 31-50 y 50+).

\begin{lstlisting}[language=python]
# Convertir fechas a datetime
datos['FECHA_NACIMIENTO'] = pd.to_datetime(datos['FECHA_NACIMIENTO'], format='%Y-%m-%d', errors='coerce')
datos['FECHA_DESAPARICION'] = pd.to_datetime(datos['FECHA_DESAPARICION'], format='%Y-%m-%d %H:%M:%S', errors='coerce')
datos['FECHA_REGISTRO'] = pd.to_datetime(datos['FECHA_REGISTRO'], format='%Y-%m-%d %H:%M:%S', errors='coerce')

# Discretizar EDAD
datos['EDAD'] = (pd.Timestamp.now() - datos['FECHA_NACIMIENTO']).dt.days / 365.25
datos['GRUPO_EDAD'] = pd.cut(datos['EDAD'], bins=[0, 12, 18, 30, 50, 100],
                              labels=['0-12', '13-17', '18-30', '31-50', '50+'])

# Discretizar FECHA_DESAPARICION
datos['MES_DESAPARICION'] = datos['FECHA_DESAPARICION'].dt.month
datos['ANO_DESAPARICION'] = datos['FECHA_DESAPARICION'].dt.year
\end{lstlisting}

Luego, removimos las columnas originales de donde viene esta información para poder trabajar con lo que ya se tiene discretizado y aquellas que no contienen información que nos sea de utilidad para ocupar los algoritmos.

\begin{lstlisting}[language=python]
columnas_a_remover = [
    'ID_VICTIMA',
    'CVE_ENT',
    'CVE_MUN',
    'FECHA_NACIMIENTO',
    'FECHA_DESAPARICION',
    'FECHA_REGISTRO',
]

datos_limpio = datos.drop(columns=columnas_a_remover, errors='ignore')
datos_limpio = datos_limpio[~(datos_limpio == 'CONFIDENCIAL').all(axis=1)]
\end{lstlisting}

Algo que notamos con una revisión manual del dataset es que existen muchas entradas con un valor 'CONFIDENCIAL'. Decidimos
dejar estos valores para poder llevar a una interpretación de los mismos posteriormente.

Este dataset modificado será usado tanto en nuestra implementación de Apriori como usando la de \texttt{mlxtend}.

\section{Implementación Propia de Apriori}

Para nuestra implementación lo que realizamos fue crear una clase (y posteriormente un objeto) \texttt{Apriori} que incluyera los métodos necesarios para llevar a la ejecución del algortimo (cálculo de lift, de confianza, etc.). En general, usamos \texttt{frozensets} para tener conjuntos inmutables y qué de esta manera fuera más fácil eliminar duplicados, trabajar con subconjuntos iguales sin orden (\{A, B\} siendo lo mismo que \{B, A\}) y que nuestra ejecución fuera más optima a largo plazo.

Decidimos incluir dos cosas extra dentro del mismo objeto Apriori. La primera fue la funcionalidad para crear las reglas de asociación. Esto para que, al momento de leer el código, fuera más directo entender de donde venían dichas reglas (pues se invocaría haciendo \texttt{Apriori.genera\_reglas\_asociacion()}). La segunda fue la capacidad de realizar el cálculo de soporte individual para poder obtener el valor más cercano a usar en el algoritmo (resultando en este caso $0.05$).

Como nota aparte, por nuestras pruebas manuales, decidimos ocupar un valor de confianza mínima de $0.3$ en el algoritmo pues los resultados tenían más sentido a la realidad.

El código de esta implementación se ve puede leer \href{https://github.com/Yoba7-7/AyMdD_2026-2/blob/main/Tarea%203/scripts/apriori.py}{en este archivo en GitHub}.

Mientras que el de la ejecución se ve a continuación:

\begin{lstlisting}[language=python]
apriori_manual = Apriori(soporte_min=0.05, confianza_min=0.3)
apriori_manual.carga_transacciones(datos_limpio, columnas=datos_limpio.columns)
itemsets_manual = apriori_manual.apriori()
reglas_asociacion = apriori_manual.genera_reglas_asociacion()
df_reglas_manual = apriori_manual.obtiene_reglas_dataframe()
# Guardado como .csv para procesamiento posterior
df_reglas_manual.to_csv('datos_apriori_manual.csv', encoding='utf-8')
\end{lstlisting}

Esta ejecución, en promedio, tomó 3.5 minutos en cada vuelta. Por que en nuestro caso, no fue la implementación más optima.

\subsection{Resultados Obtenidos}

A continuación se incluye una muestra de las 68 reglas de asociación obtenidas. El archivo completo de ellas se encuentra
disponible \href{https://github.com/Yoba7-7/AyMdD_2026-2/blob/main/Tarea%203/datasets/datos_apriori_manual.csv}{en este archivo en GitHub}.

\begin{lstlisting}
,antecedente,consecuente,soporte,confianza,lift
0,frozenset({'ENTIDAD=TAMAULIPAS'}),frozenset({'ORIGEN_REPORTE=PROCURADURIA GENERAL DE JUSTICIA DEL ESTADO DE TAMAULIPAS'}),0.08967263438571332,0.8925066904549509,9.887881114186348
1,frozenset({'ORIGEN_REPORTE=PROCURADURIA GENERAL DE JUSTICIA DEL ESTADO DE TAMAULIPAS'}),frozenset({'ENTIDAD=TAMAULIPAS'}),0.08967263438571332,0.9934629706247414,9.887881114186348
...
63,frozenset({'GRUPO_EDAD=18-30'}),frozenset({'SEXO=HOMBRE'}),0.06344155892655746,0.6781095321730799,1.403984333870305
64,frozenset({'GRUPO_EDAD=18-30'}),"frozenset({'SEXO=HOMBRE', 'ESTATUS_VICTIMA=DESAPARECIDA'})",0.059595031631151645,0.6369950502953856,1.4024890034352622
65,frozenset({'ENTIDAD=ESTADO DE MEXICO'}),frozenset({'ESTATUS_VICTIMA=DESAPARECIDA'}),0.07751312674120714,0.7235080870050195,1.2126579192153457
66,frozenset({'ENTIDAD=TAMAULIPAS'}),frozenset({'ESTATUS_VICTIMA=DESAPARECIDA'}),0.06084235213276868,0.6055605114481118,1.014968267751447
67,frozenset({'ORIGEN_REPORTE=PROCURADURIA GENERAL DE JUSTICIA DEL ESTADO DE TAMAULIPAS'}),frozenset({'ESTATUS_VICTIMA=DESAPARECIDA'}),0.05429952123805896,0.6015721969383533,1.0082835308957736
68,"frozenset({'ENTIDAD=TAMAULIPAS', 'ORIGEN_REPORTE=PROCURADURIA GENERAL DE JUSTICIA DEL ESTADO DE TAMAULIPAS'})",frozenset({'ESTATUS_VICTIMA=DESAPARECIDA'}),0.05393354097111743,0.6014492753623188,1.0080775044182568
\end{lstlisting}


Sin embargo, era necesario hacer una primera selección de las mismas, ya que algunas tienen un lift $<5$, lo cual nos lleva
a considerarlas como reglas espurias. Esto nos hace eliminar las primeras 9 reglas resultantes, con una distribución de lift
que se muestra como sigue:

\begin{figure}[H]
    \centering
    \includegraphics[width=0.7\linewidth]{imagenes/grafico_apriori_reglas_manual.png}
    \label{fig:scatter-aprori-manual}
\end{figure}

Esta pre-selección de 59 reglas será usada posteriormente para comparar resultados con los obtenidos en la siguiente sección.

\section{Usando Apriori de \texttt{mlxtend}}

Para poder trabajar con \texttt{mlxtend} debemos hacer algunas modificaciones a una copia de nuestro dataset limpio (el que usamos anteriormente para ocupar nuestra implementación de Apriori). El dataframe que ocupa \texttt{mlxtend.aprori} debe ser un dataframe booleano o \textit{onehot}. Esto es, un dataframe en donde se de convertir cada valor posible de una columna en una columna booleana. Por ejemplo, si tenemos que los valores en la columna 'SEXO' son HOMBRE, MUJER, etc, entonces se debe detener un columna SEXO=HOMBRE con valores True o False, otra con SEXO=MUJER, etc, etc.

El código de esta implementación se ve de la siguiente manera:

\begin{lstlisting}[language=python]
# Para trabajar con algo que entienda el apriori de mlxtend, se debe convertir a un dataframe booleano/onehot.
# Esto es, convertir cada valor en una columna booleana.

datos_para_mlxtend = datos_limpio.copy()

# En lugar de insertar columna por columna, crear una lista de columnas
columnas_onehot = {}

for columna in datos_para_mlxtend.columns:
    valores_unicos = datos_para_mlxtend[columna].dropna().unique()

    for valor in valores_unicos:
        nombre_col = f"{columna}={valor}"
        columnas_onehot[nombre_col] = (datos_para_mlxtend[columna] == valor).astype(bool)

datos_onehot = pd.DataFrame(columnas_onehot)

itemsets_frequentes_mlxtend = apriori(datos_onehot, min_support=0.05, use_colnames=True)
reglas_mlxtend = association_rules(itemsets_frequentes_mlxtend, metric='confidence', min_threshold=0.3)
reglas_mlxtend = reglas_mlxtend[['antecedents', 'consequents', 'support', 'confidence', 'lift']]
reglas_mlxtend.to_csv('datos_mlxtend.csv', encoding='utf-8')
\end{lstlisting}

Esta ejecución en promedio tomó 35 segundos segundos en cada vuelta, lo que lo hace considerablemente más rápido y mejor optimizado que nuestra implementación.

\subsection{Resultados Obtenidos}

Sin embargo, aunque las reglas de asociación obtenidas con este método estuvieron en un orden distinto, obtuvimos las mismas 68 reglas de asociación que con nuestra implementación manual. He aquí una muestra de los resultados:

\begin{lstlisting}
,antecedents,consequents,support,confidence,lift
0,frozenset({'ORIGEN_REPORTE=PROCURADURIA GENERAL DE JUSTICIA DEL ESTADO DE TAMAULIPAS'}),frozenset({'ESTATUS_VICTIMA=DESAPARECIDA'}),0.05429952123805896,0.6015721969383533,1.0082835308957736
1,frozenset({'ENTIDAD=TAMAULIPAS'}),frozenset({'ORIGEN_REPORTE=PROCURADURIA GENERAL DE JUSTICIA DEL ESTADO DE TAMAULIPAS'}),0.08967263438571332,0.8925066904549509,9.887881114186348
2,frozenset({'ORIGEN_REPORTE=PROCURADURIA GENERAL DE JUSTICIA DEL ESTADO DE TAMAULIPAS'}),frozenset({'ENTIDAD=TAMAULIPAS'}),0.08967263438571332,0.9934629706247414,9.887881114186348
...
66,"frozenset({'ENTIDAD=JALISCO', 'SEXO=CONFIDENCIAL'})","frozenset({'ESTATUS_VICTIMA=CONFIDENCIAL', 'MUNICIPIO=CONFIDENCIAL'})",0.07468985039622966,1.0,2.7241042544100593
67,"frozenset({'ENTIDAD=JALISCO', 'MUNICIPIO=CONFIDENCIAL'})","frozenset({'ESTATUS_VICTIMA=CONFIDENCIAL', 'SEXO=CONFIDENCIAL'})",0.07468985039622966,1.0,2.7241042544100593
68,frozenset({'ENTIDAD=JALISCO'}),"frozenset({'ESTATUS_VICTIMA=CONFIDENCIAL', 'MUNICIPIO=CONFIDENCIAL', 'SEXO=CONFIDENCIAL'})",0.07468985039622966,0.6751738572682466,1.8392439770508808
\end{lstlisting}


El archivo completo de ellas se encuentra disponible \href{https://github.com/Yoba7-7/AyMdD_2026-2/blob/main/Tarea%203/datasets/datos_mlxtend_manual.csv}{en este archivo en GitHub}.

De igual manera, se hace una preselección de reglas con lift$<5$ obteniendo, de igual manera, 59 reglas de este conjunto:

\begin{figure}[H]
    \centering
    \includegraphics[width=0.7\linewidth]{imagenes/grafico_apriori_reglas_mlxtend.png}
    \label{fig:scatter-aprori-mlxtend}
\end{figure}

En consecuente, en lugar de hacer una comparativa de las reglas obtenidas, procedemos con la selección de las mismas para llegar a la interpretación de los resultados.

\section{Interpretación de las Reglas de Asociación}

\subsection{Rangos de edad}

Tenemos reglas que nos indican que los rangos de edad más comunes entre las víctimas desaparecidas.

\begin{table}[H]
\centering
\small
\begin{tabular}{cccccc}
\hline
ID & Antecedente & Consecuente & Soporte & Confianza & Lift \\
\hline
24 & \texttt{GRUPO\_EDAD=$50+$} & \texttt{ESTATUS\_VICTIMA=DESAPARECIDA} & 0.0840 & 0.9314 & 1.5611 \\
25 & \texttt{GRUPO\_EDAD=31-50} & \texttt{ESTATUS\_VICTIMA=DESAPARECIDA} & 0.1881 & 0.3154 & 1.5723 \\
26 & \texttt{ESTATUS\_VICTIMA=DESAPARECIDA} & \texttt{GRUPO\_EDAD=31-50} & 0.1882 & 0.3154 & 1.5723 \\
27 & \texttt{GRUPO\_EDAD=18-30} & \texttt{ESTATUS\_VICTIMA=DESAPARECIDA} & 0.0874 & 0.9348 & 0.9348 \\
\hline
\end{tabular}
\caption{Regla directa para rango de edad condicionado a estatus desaparecida. Valores redondeados.}
\label{tab:edad_dado_desaparecida}
\end{table}

Estas reglas sugieren que el rango de edad promedio de las víctimas desaparecidas es entre 31 y 50 años (con valores de
confianza de 0.3154). Esto puede indicar que las personas en este rango de edad son más vulnerables a la desaparición, o
que hay un sesgo en el registro hacia estos grupos demográficos. Sin embargo, también es importante considerar que el rango
de edad de 18-30 años tiene una confianza más alta (0.9348) para el estatus de desaparecida, lo que podría indicar que aunque
hay menos casos en este rango, es más probable que sean reportados como desaparecidos.

\subsection{Perfil Demográfico}

Las siguientes reglas sugieren un perfil demográfico: conforme aumenta la edad, se observa una mayor sobrerrepresentación de hombres en el registro:

\begin{table}[H]
\centering
\small
\begin{tabular}{cccccc}
\hline
ID & Antecedente & Consecuente & Soporte & Confianza & Lift \\
\hline
14 & \texttt{GRUPO\_EDAD=50+} & \texttt{SEXO=HOMBRE} & 0.0790 & 0.8756 & 1.813 \\
15 & \texttt{GRUPO\_EDAD=31-50} & \texttt{SEXO=HOMBRE} & 0.1671 & 0.8330 & 1.725 \\
17 & \texttt{GRUPO\_EDAD=18-30} & \texttt{SEXO=HOMBRE} & 0.0634 & 0.6781 & 1.404 \\
53 & \texttt{SEXO=HOMBRE, ESTATUS\_VICTIMA=DESAPARECIDA} & \texttt{GRUPO\_EDAD=31-50} & 0.1567 & 0.3449 & 1.719 \\
\hline
\end{tabular}
\caption{Reglas de asociación para perfil demográfico (edad y sexo). Valores redondeados.}
\label{tab:perfil_demografico}
\end{table}

Aquí, los lifts crecientes ($<1$) por grupo sugieren un sesgo de reporte/captura por edad y sexo.

\subsection{Confidencialidad en Bloque}

Estas reglas muestran un patrón sistémico: cuando un campo es confidencial, los demás también lo son, con una confianza
de 1.0, en una proporción muy grande del dataset.

\begin{table}[htbp]
\centering
\small
\begin{tabular}{cccccc}
\hline
ID & Antecedente & Consecuente & Support & Conf. & Lift \\
\hline
7  & \texttt{ESTATUS\_VICTIMA=CONFIDENCIAL} & \texttt{SEXO=CONFIDENCIAL} & 0.3671 & 1.0000 & 2.724 \\
10 & \texttt{SEXO=CONFIDENCIAL} & \texttt{MUNICIPIO=CONFIDENCIAL} & 0.3671 & 1.0000 & 2.724 \\
20 & \texttt{ESTATUS\_VICTIMA=CONFIDENCIAL} & \texttt{MUNICIPIO=CONFIDENCIAL} & 0.3671 & 1.0000 & 2.724 \\
\hline
\end{tabular}
\caption{Sesgo de registro: confidencialidad en bloque (censura/datos faltantes estructurados). Valores redondeados.}
\label{tab:confidencialidad_bloque}
\end{table}

Esto nos lleva a esta conclusión: existe un mecanismo de censura sistémico de la información. Se realiza el registro
de la persona desaparecida, pero no se deciden hacer públicos ninguno de los datos personales. Esto puede deberse a
políticas de privacidad (sin explicar cuales son), temor a represalias, o falta de información capturada en el momento
del registro.

\subsection{Sesgo por entidad — caso Jalisco}

Estas reglas son un triste ejemplo de lo anterior. Si el registro de desaparición se hace en Jalisco, es muy probable
que el resto de los campos terminen siendo confidenciales, lo que impide cualquier análisis de perfil o tendencia en esta entidad.

\begin{table}[H]
\centering
\small
\begin{tabular}{cccccc}
\hline
ID & Antecedente & Consecuente & Soporte & Confianza & Lift \\
\hline
9  & \texttt{ENTIDAD=JALISCO} & \texttt{SEXO=CONFIDENCIAL} & 0.0747 & 0.6752 & 1.839 \\
19 & \texttt{ENTIDAD=JALISCO} & \texttt{ESTATUS\_VICTIMA=CONFIDENCIAL} & 0.0747 & 0.6752 & 1.839 \\
30 & \texttt{ENTIDAD=JALISCO} & \texttt{MUNICIPIO=CONFIDENCIAL} & 0.0747 & 0.6752 & 1.839 \\
36 & \texttt{ESTATUS\_VICTIMA=CONFIDENCIAL, ENTIDAD=JALISCO} & \texttt{SEXO=CONFIDENCIAL} & 0.0747 & 1.0000 & 2.724 \\
45 & \texttt{ENTIDAD=JALISCO', 'SEXO=CONFIDENCIAL} & \texttt{MUNICIPIO=CONFIDENCIAL} & 0.0747 & 1.0000 & 2.724 \\
46 & \texttt{ENTIDAD=JALISCO', 'MUNICIPIO=CONFIDENCIAL} & \texttt{SEXO=CONFIDENCIAL} & 0.0747 & 1.0000 & 2.724 \\
\hline
\end{tabular}
\caption{Sesgo por entidad: patrones de confidencialidad asociados a Jalisco. Valores redondeados.}
\label{tab:jalisco_confidencialidad}
\end{table}

\subsection{Sesgo Temporal}

En base a estas dos reglas podemos inferir que los registros de 2024 en adelante
tienen un rezago de actualización, ya que no existe existe un cambio en el estado
de la víctima. Esto se puede deber a que el proceso de investigación y actualización
de casos es lento, o a que los casos más recientes aún no han sido actualizados en el
sistema.

\begin{table}[H]
\centering
\small
\begin{tabular}{cccccc}
\hline
ID & Antecedente & Consecuente & Soporte & Confianza & Lift \\
\hline
28 & \texttt{AÑO\_DESAPARICION=2025.0} & \texttt{ESTATUS\_VICTIMA=DESAPARECIDA} & 0.0512 & 0.9985 & 1.674 \\
29 & \texttt{AÑO\_DESAPARICION=2024.0} & \texttt{ESTATUS\_VICTIMA=DESAPARECIDA} & 0.0598 & 0.9960 & 1.669 \\
\end{tabular}
\caption{Sesgo temporal: asociación entre año reciente y estatus ``DESAPARECIDA'' (posible rezago de actualización). Valores redondeados.}
\label{tab:sesgo_temporal}
\end{table}

Los archivos con los cuales se hicieron estas interpretaciones se pueden encontrar
aqui:
\begin{itemize}
    \item \href{https://github.com/Yoba7-7/AyMdD_2026-2/tree/main/Tarea%203/datasets#:~:text=datos_manual_final.csv}{Reglas Finales con nuestra
    implementación de Apriori.}
    \item \href{https://github.com/Yoba7-7/AyMdD_2026-2/blob/main/Tarea%203/datasets/datos_mlxtend_final.csv}{Reglas Finales con la implementación de mlxtend.}
\end{itemize}